\subsection{Mimo}

% Final
První analyzovanou aplikací je aplikace Mimo\footnote{Dostupné na: \href{https://mimo.org/}{https://mimo.org/}}, která se nejvíce blíží vizi tohoto projektu a bude tak přímou konkurencí vznikající aplikace. Mimo se zaměřuje zejména na výuku moderních programovacích jazyků a frameworků. V době psaní této práce Mimo nabízí 9 kurzů, 4 tzv. kariérní cesty a má napříč platformami přes 40 milionů registrovaných uživatelů \cite{Mimo}. Mimo je k dispozici jako webová aplikace a dále také jako mobilní aplikace na platformách Android a iOS. Pro analýzu použiji primárně aplikaci na platformě Android.

% Final
\subsubsection{Vzhled a uživatelská zkušenost}

% Odstavec: Obecný dojem (komplexita/přehlednost UI, paleta barev, typografie)
% Final
Mimo se od dalších konkurentů liší zejména svým vzhledem uživatelského rozhraní a stylem výuky. Uživatelské rozhraní je velice minimalistické a přehledné. Místy může aplikace působit až příliš jednoduchým dojmem. Aplikace využívá fialovou monochromatickou paletu barev a celé rozhraní využívá spíše tmavší barvy. V rámci typografie Mimo využívá pro některé nadpisy neproporcionální písmo. To lze pozorovat zejména v mobilní aplikaci a na webových stránkách. Pro textové prvky pak využívá bezpatkové písmo Aeonik.

% Odstavec: Rozložení prvků na obrazovce (lekce v kurzu, učení)
% Final
Rozložení prvků na obrazovce je přehledné. Na hlavní stránce kurzu jsou umístěny jednotlivé lekce, které jsou vyobrazeny ve formě tlačítek na čáře připomínající cestu kurzem. Až po kliknutí na danou lekci se zobrazí její název a možnost spuštění dané lekce. Toto vyobrazení lekcí může na uživatele působit vizuálně poutavým dojmem, nicméně může být obtížné se pomocí něj orientovat v kurzu, protože o lekcích nejsou na první pohled zobrazeny žádné podrobnější informace.

% Odstavec: Navigace v aplikaci
% Final
Pro navigaci v mobilní aplikaci používá Mimo velice jednoduchou spodní navigační lištu s pouze pěti prvky. Díky tomu je navigace v aplikaci velice snadná a rychlá.

% Odstavec: Srovnání webového a mobilního prostředí
% Final
Webová verze aplikace Mimo je velice podobná mobilní verzi. Rozdíl je především v tom, že ve webové aplikaci jsou lekce vyobrazeny ve formě seznamu. Pro navigaci je použita horní navigační lišta a pro navigaci mezi kapitolami kurzu pak boční navigační lišta.

% Odstavec: Animace a efekty a mikro interakce
% Final
V aplikaci jsou použity animace a efekty velice zřídka. Některé prvky uživatelského rozhraní jsou animovány při jejich zobrazení. Při učení jsou také použity animace pro přechody mezi jednotlivými cvičeními.

% Final
\subsubsection{Způsoby učení}

% Struktura kurzu (lineární, tematické moduly, career paths, kapitoly atd.)
% Final
Jednotlivé kurzy jsou uspořádány do sekcí, které pak obsahují jednotlivé lekce. Lekce jsou uspořádány lineárně a uživatel se je musí učit v daném pořadí. Kromě kurzů Mimo nabízí tzv. kariérní cesty, které mají podobnou strukturu jako kurzy, nicméně většinou obsahují více témat.

% Jak se spustí učení a jak vypadá study session
% Final
Učení lze spustit kliknutím na danou lekci. Při učení se uživateli zobrazuje sekvence stránek na nichž se nachází buď vysvětlení dané látky nebo různá cvičení. Při vyplňování cvičení je typicky uživateli zobrazena otázka, na kterou musí odpovědět například vyplněním textového pole v kódu. Po vyplnění odpovědi se uživateli zobrazí, zda byla odpověď správná, a u cvičení se spustitelným kódem se uživateli zobrazí výstup takového kódu. V kurzu SQL je tak uživateli například po vyplnění odpovědi zobrazena tabulka s daty, která by byla výsledkem daného SQL dotazu. Po vyplnění odpovědi ve webové aplikaci má uživatel také možnost si nechat odpověď vysvětlit v chatovém okně pomocí chatbota.

% Druhy cvičení
% Final
V aplikaci se nachází několik typů cvičení, které lze rozdělit do následujících kategorií.

% Final
\paragraph{Otázka s více odpovědmi} Jedním z nejčastějších typů cvičení je otázka, která typicky obsahuje ukázku kódu, a výběr z několika odpovědí. Uživatel má za cíl zvolit jednu správnou odpověď.

% Final
\paragraph{Doplňování do kódu} Dalším častým typem cvičení je doplňování částí kódu. Aplikace uživateli v otázce zobrazí ukázku kódu, do které je třeba doplnit typicky několik slov nebo znaků. Doplňování je v některých případech usnadněno tím, že může uživatel slova a znaky vybírat z předdefinovaného seznamu. To může značně zrychlit procházení těchto cvičení, protože uživatel nemusí vše psát ručně.

% Final
Dále aplikace ve cvičeních využívá řadu komponentů pro vysvětlení a lepší pochopení dané látky. Příkladem toho jsou následující komponenty.

% Final
\paragraph{Tabulky} V některých kurzech, zejména v kurzech vyučující témata týkající se databází, jsou využívány tabulky s daty. Ty se zobrazují jednak jako součástí zadání cvičení a jednak jako výstup po vyplnění daného cvičení.

% Final
\paragraph{Výstup kódu v terminálu} Po vyplnění cvičení, které v zadání obsahuje spustitelný kód, může být uživateli zobrazen komponent připomínající terminál s výstupem takového kódu. Tyto komponenty jsou typicky zobrazovány například u kurzů, které vyučují konkrétní programovací jazyky.

% Final
\paragraph{Zobrazení webové stránky} Po vyplnění cvičení, které v zadání obsahuje kód a zaměřuje se na webové technologie, je typicky zobrazena ukázka webové stránky, která reprezentuje výstup daného kódu. Tento typ cvičení je využívaný zejména u výuky jazyků, jako například HTML, CSS nebo JavaScript. Dále je tento typ cvičení také využíván pro výuku frontend frameworků.

% Final
\subsubsection{Způsoby motivace uživatele}

% Jaké formy gamifikace aplikace používá
% Final
Mimo používá různé způsoby motivace uživatelů k učení. Jedním z těchto způsobů je gamifikace, kterou lze definovat jako použití prvků herního designu v mimoherních kontextech \cite{10.1145/2181037.2181040} (překlad autora). V aplikaci Mimo jsem identifikoval následující formy gamifikace.

% Final
\paragraph{Virtuální měna} V aplikaci je možné za plnění cvičení získávat virtuální měnu, za kterou je následně možné si kupovat různé položky v obchodě aplikace.

% Final
\paragraph{Streak} Aplikace zaznamenává, kolik dní v řadě uživatel splnil alespoň jednu lekci. Tato řada je nazývána streak. Pokud uživatel daný den nesplnil žádnou lekci, je mu jeho streak resetován. V obchodě aplikace si může uživatel zakoupit několik položek, které mu pomohou zachovat streak i přes to, že se daný den zapomněl učit. Uživatel si tak může preventivně zakoupit položku, která mu zajistí, že se streak nezruší, pokud se zapomene některý den učit. Svůj streak může uživatel také zpětně napravit zakoupením příslušné položky.

% Final
\paragraph{Streak výzva} V obchodě aplikace může uživatel za virtuální měnu zakoupit tzv. streak výzvu. Pokud uživatel splní výzvu, která spočívá v učení se každý den po dobu 7 dní, dostane jako odměnu určité množství virtuální měny.

% Final
\paragraph{Srdíčka} Uživatel má k dispozici 5 srdíček, o která přichází, pokud neodpoví správně na danou otázku v průběhu cvičení. Srdíčka se obnovují po určitém čase a uživatel je také může zakoupit za virtuální měnu.

% Final
\paragraph{Body XP} Za plnění jednotlivých lekcí v aplikaci uživatel může získat tzv. body XP, které reprezentují pokrok uživatele v rámci aplikace. Tyto body jsou následně zobrazeny na jeho profilu a ovlivňují jeho pozici v lize učení.

% Final
\paragraph{Liga učení} Uživatelé jsou každý týden rozřazeni do ligy, neboli skupiny uživatelů, na základě toho, kolik bodů XP získali v předchozím týdnu. Na profilu uživatele se pak zobrazuje, ve které lize se nachází.

% Final
\paragraph{Systém přátel} Mimo umožňuje uživatelům si přidávat ostatní uživatele do přátel a sledovat tak jejich pokrok v aplikaci.

% Final
\paragraph{Ikona aplikace} Mimo umožňuje uživateli si za virtuální měnu zakoupit a změnit ikonu aplikace.

% Final
\subsubsection{Placená verze}

% Free tier
% Final
Mimo nabízí uživatelům zdarma pouze prvních několik lekcí každého kurzu. Při učení se uživatelům zobrazují reklamy a učení je omezeno funkcí srdíček.

% Premium tier
% Final
Uživatelé, kteří si zaplatí premium účet, se mohou všechny kurzy a kariérní cesty učit bez omezení. Při učení se jim nezobrazují reklamy a nejsou omezeni funkcí srdíček. Uživatelé s premium účtem si navíc mohou po dokončení kurzu vygenerovat a stáhnout certifikát o dokončení daného kurzu. 

% Free tier
% Final
Aplikace nabízí sedmidenní zkušební verzi premium účtu zdarma na webových stránkách Mimo. Zkušební verzi lze také získat pozváním přátel do aplikace.

% Final
\subsubsection{Silné a slabé stránky}

% Silné stránky
% Final
Silnou stránkou aplikace Mimo je její jednoduchost používání a krátké lekce. Uživatel se tak může snadno učit v průběhu dne čistě s využitím mobilu, což může být praktické zejména pro uživatele, kteří mají omezený čas na učení.

% Final
Další silnou stránkou je použitá gamifikace, která může vést ke zvýšené motivaci uživatelů používat aplikaci pravidelně.

% Slabé stránky
% Final
Možnou nevýhodou minimalistického vzhledu aplikace může být pocit, který uživatelské rozhraní svým minimalismem vyvolává. Aplikace může místy působit až příliš jednoduše a vyvolat tak dojem, že je příliš jednoduchá na to, aby vyučovala komplexní a pokročilá témata, kterými jsou například programování a softwarové inženýrství. To by mohlo některé uživatele odradit, zejména ty, kteří potřebují daná témata znát na profesionální úrovni.

% Final
Další potenciální nevýhodou je samotná navigace v kurzu. Lekce kurzu jsou vyobrazeny vizuálně poutavým způsobem, nicméně najít konkrétní lekci kurzu může být pro uživatele velice obtížné, zejména v mobilní verzi aplikace.

% Final
\subsubsection{Další poznatky}

% Zajímavé funkce
% Final
Mimo kromě základních funkcí poskytuje také tzv. pískoviště, kde může uživatel vytvářet a spouštět svůj kód a vidět jeho výstup. V rámci této funkce jsou k dispozici pouze jazyky HTML, CSS a JavaScript a dále framework React.

% Final
Aplikace také poskytuje funkci, v rámci které může uživatel vytvořit vlastní webovou aplikaci přímo uvnitř Mimo. V rámci této funkce uživatel zadá umělé inteligenci, co chce vytvořit. Mimo pak uživateli nastaví virtuální prostředí se systémem souborů, kde umělá inteligence celý projekt připraví. Uživatel pak může tento projekt upravovat a spouštět přímo v Mimo. Uživatel si může v rámci této funkce zobrazit danou webovou stránku, soubory projektu a také jednoduché zobrazení tabulek v databázi. Celý projekt pak uživatel může publikovat na veřejně dostupné webové adrese. Tato funkce je pro uživatele bez premium účtu značně omezená.

% (Zajímavé detaily nebo cokoliv co mě napadne)
